\section{Matrici}
\emph{[Rif. Lezioni 2 e 3]}

\subsection{Definizioni Fondamentali}
Per studiare i sistemi lineari, è fondamentale introdurre il concetto di matrice.

\begin{definitionbox}
\textbf{Matrice:} Una matrice $A$ è una tabella rettangolare di numeri disposti su $m$ righe e $n$ colonne.
\begin{itemize}
    \item \textbf{Matrice Associata ($A$):} Tabella $m \times n$ dei coefficienti. $A = (a_{ij})_{i=1..m, j=1..n}$.
    \item \textbf{Vettore Termini Noti ($b$):} Matrice colonna $m \times 1$.
    \item \textbf{Matrice Completa ($A|b$):} Matrice $m \times (n+1)$ ottenuta affiancando $b$ ad $A$.
\end{itemize}
\end{definitionbox}

\subsection{Tipologie Speciali di Matrici}
Una matrice si dice \textbf{Quadrata} di \textbf{ORDINE} $n$ se ha lo stesso numero di righe e colonne ($m=n$).

\begin{concept}{Matrici Triangolari e Diagonali}
\begin{itemize}
    \item \textbf{Diagonale Principale:} Insieme degli elementi $a_{ii}$ ($a_{11}, a_{22}, \dots$).
    \item \textbf{Matrice Diagonale:} Tutti gli elementi fuori dalla diagonale principale sono nulli ($a_{ij} = 0$ per $i \neq j$).
    $$ A = \begin{pmatrix} d_1 & 0 & 0 \\ 0 & d_2 & 0 \\ 0 & 0 & d_3 \end{pmatrix} $$
    \item \textbf{Matrice Triangolare Superiore:} Tutti gli elementi \textbf{SOTTO} la diagonale sono nulli ($a_{ij}=0 \ \forall i > j$).
    \item \textbf{Matrice Triangolare Inferiore:} Tutti gli elementi \textbf{SOPRA} la diagonale sono nulli ($a_{ij}=0 \ \forall i < j$).
\end{itemize}
\end{concept}

\begin{example}
Esempio di matrice triangolare superiore:
$$ A = \begin{pmatrix} 1 & 3 & 5 \\ 0 & 2 & 4 \\ 0 & 0 & 6 \end{pmatrix} $$
\end{example}
